% Stand-in build: the standard article class, kept so that the manuscript can be
% compiled and read without the MDPI class. The submission version is main.tex,
% which uses Definitions/mdpi.cls. Both pull the same sections/ files, so any
% edit to a section shows up in both.
%
% Build: "D:/environment/tools/tectonic/tectonic.exe" main_article.tex
\documentclass[11pt,a4paper]{article}

\usepackage{amsmath}
\usepackage{amssymb}
\usepackage{amsthm}
\usepackage{algorithm}
\usepackage{algpseudocode}
\usepackage{booktabs}
\usepackage{array}
\usepackage{graphicx}
% The four figures are drawn in LaTeX from numbers that are already in the
% tables or in the cited precheck files; no external image is included.
\usepackage{tikz}
\usetikzlibrary{decorations.pathmorphing,positioning,arrows.meta}
\usepackage{pgfplots}
\pgfplotsset{compat=1.18}
\usepackage[hidelinks]{hyperref}

% Three wide tables are set across the full page with the MDPI template's
% adjustwidth/\extralength mechanism, where \extralength is the width of the
% journal's left margin column. The article class has no such column, so the
% two lengths are given here: 70 pt of the 125 pt left margin is lent to those
% tables, which is what keeps them inside the measure now that the page-saving
% \scriptsize has been removed. Setting \extralength to 0 pt puts the tables
% back inside the text block.
\usepackage{changepage}
\newlength{\extralength}
\setlength{\extralength}{70pt}
\newlength{\fulllength}
\setlength{\fulllength}{\textwidth}
\addtolength{\fulllength}{\extralength}

% unsrt prints the note field. The verification records that used to sit there,
% including the Chinese ones that rendered as blanks under a Latin-only font,
% now live in an x-verified field, which no bst reads; only the short
% bibliographic notes (Accessed ..., Preprint arXiv:..., Author Correction) print.
\bibliographystyle{unsrt}

% Theorem environments are declared here and only here. sections/06_theory.tex
% re-declares them only when they are still undefined, so this preamble wins.
% main.tex instead maps these names onto the environments mdpi.cls provides.
\theoremstyle{plain}
\newtheorem{theorem}{Theorem}[section]
\newtheorem{lemma}[theorem]{Lemma}
\newtheorem{corollary}[theorem]{Corollary}
\newtheorem{proposition}[theorem]{Proposition}
\theoremstyle{definition}
\newtheorem{definition}[theorem]{Definition}
\theoremstyle{remark}
\newtheorem{remark}[theorem]{Remark}

% The sections wrap align in linenomath, which mdpi.cls (through lineno) needs for line
% numbers inside multi-line displays; the article build has no line numbers, so the
% environment is empty here.
\newenvironment{linenomath}{}{}

% Marks a figure that comes from development data rather than from the sealed
% semesters. It expands to its argument; the submission version can redefine it.
\newcommand{\devnum}[1]{#1}
% Figures from the sealed terms, opened once after the protocol freeze.
\newcommand{\sealednum}[1]{#1}

% Pulls in a section if its file exists, and otherwise leaves a numbered
% placeholder so that section numbers and cross-references stay correct.
% Adjust the file name here if a section is committed under a different one.
\newcommand{\papersection}[3]{%
  \IfFileExists{sections/#1.tex}%
    {\input{sections/#1}}%
    {\section{#2}\label{#3}%
     \textit{Not yet in the repository: the file
     \texttt{\detokenize{sections/#1.tex}} was not found.}}%
}

\title{Prediction-Driven Non-Preemptive Scheduling with Bounded Overtaking
for Shared Execution Services under Deadline-Driven Bursty Load}

\author{[Authors]}

\date{}

\begin{document}

\maketitle

\begin{abstract}
A shared execution service runs submitted jobs on a fixed pool of servers,
without preemption, under a per-job time limit $L$. Reordering by predicted cost
shortens most waits and can leave one job far behind first-come first-served. We
prove a conservation law. For every non-preemptive work-conserving policy on $k$
servers, every arrival sequence and every job, the delay relative to first-come
first-served equals the net overtaken work divided by $k$, to within
$(2-2/k)L$, an exact supremum, approached and never attained. Bounding overtaken
work is therefore sufficient and necessary for a per-job guarantee; counting
overtakes is neither. A wrapper on that budget bounds every job's wait against
first-come first-served, additively and multiplicatively, under arbitrary
arrivals and predictions; its additive constant $(3-2/k)L$ is likewise an exact
supremum, and the operator sets it from a promise in seconds. The score underneath should estimate expected
cost, not a point prediction on a log scale. We evaluate on two judging
platforms, a serverless invocation trace, two continuous-integration pools whose
recorded waits validate the simulator, and a compute farm. The method needs a
time limit comparable to the waits at issue and a heavy-tailed cost
distribution; we include a workload failing each.
\end{abstract}

% The MDPI class supplies \keyword{}; under the stand-in class the keywords are
% typeset by hand here.
\noindent\textit{Keywords:} non-preemptive scheduling; scheduling with
predictions; multiserver queues; bounded overtaking; per-job guarantee;
first-come first-served; runtime prediction

\bigskip

\papersection{01_introduction}{Introduction}{sec:intro}
\papersection{02_related_work}{Related Work}{sec:related}
\papersection{03_problem_model}{Problem Model}{sec:model}
\papersection{04_prediction}{Cost Prediction}{sec:prediction}
\papersection{05_scheduling}{Scheduling Algorithm}{sec:scheduling}
\papersection{06_theory}{Theory}{sec:theory}
\papersection{07_data}{Traces and Workload Characterisation}{sec:data}
\papersection{08_experiments}{Experiments}{sec:experiments}
\papersection{09_limitations}{Limitations and Threats to Validity}{sec:limits}

\IfFileExists{sections/A_proofs.tex}{\providecommand{\devnum}[1]{#1}

\appendix

\section{Proofs of the tightness results}
\label{app:proofs}

This appendix gives in full the constructions sketched in
Section~\ref{sec:theory}. All of them are exact in integer arithmetic; we
simulated every one of them and checked every intermediate quantity of every
proof job by job, including an independent audit of work conservation and server
capacity re-derived from the produced schedule.

\subsection{The workload comparison}
\label{app:strict}

\begin{proof}[Proof of Lemma~\ref{lem:gap}]
Put $D = U_A - U_B$. Arrivals add the same jump to both terms, so $D$ is
continuous, and between arrivals it changes at rate
$-(\min(N_A,k) - \min(N_B,k))$; it is therefore piecewise linear with
breakpoints among the finitely many arrival and completion instants of the two
runs.

If $U_A(u) > (k-1)L$ then more than $k-1$ jobs are present under $A$, each
carrying at most $L$, so every server of $A$ is busy and
$\min(N_A,k) = k \ge \min(N_B,k)$: $D$ does not increase at $u$. Fix $t$ and let
$S = \{u \le t : U_A(u) \le (k-1)L\}$, which is non-empty because $U_A$ vanishes
before the first arrival, and put $t_0 = \sup S$. The supremum need not be
attained, since $U_A$ jumps upwards at arrivals, so we use the left limit:
there are $u_n \uparrow t_0$ with $U_A(u_n) \le (k-1)L$, hence
$U_A(t_0^-) \le (k-1)L$ and, $D$ being continuous,
$D(t_0^-) \le U_A(t_0^-) \le (k-1)L$ because $U_B \ge 0$. On $(t_0, t]$ we have
$U_A > (k-1)L$, so $D$ does not increase there, and $D(t) \le (k-1)L$.
Exchanging $A$ and $B$ gives the other direction. That the constant cannot be
lowered is the cascade of Appendix~\ref{app:cascade}, which reaches
$(k-1)L\,(1 - ((k-1)/k)^m)$ after $m$ rounds.
\end{proof}

\begin{proof}[Proof of the strict inequality in Lemma~\ref{lem:gap}]
Keep the notation above, so $D = U_A - U_B$ is continuous and piecewise
linear with finitely many pieces. Suppose $D(t) = (k-1)L$ for some $t$ and let
$t_1$ be the smallest such $t$; it exists because $\{D = (k-1)L\}$ is closed and
$D(0) = 0 < (k-1)L$, using $L > 0$.

Let $I$ be the linear piece immediately to the left of $t_1$, which has positive
length. Its slope is strictly positive: a negative slope would put
$D > (k-1)L$ somewhere on $I$ and a zero slope would put $D = (k-1)L$ on all of
$I$, and either contradicts the minimality of $t_1$. The slope of $D$ is
$-(\min(N_A,k) - \min(N_B,k))$, so $\min(N_A,k) < \min(N_B,k)$ on $I$ and
therefore $N_A \le k-1$ there.

With at most $k-1$ jobs present and $A$ work-conserving, none of $A$'s jobs
waits, so all of them are in service; hence $U_A$ falls at rate $N_A$ on $I$,
and $U_A \le (k-1)L$ on $I$ because each of the at most $k-1$ present jobs
carries at most $L$. Taking left limits at $t_1$ and using the continuity of
$D$,
\[
  U_A(t_1^-) - U_B(t_1^-) = (k-1)L , \qquad U_A(t_1^-) \le (k-1)L ,
  \qquad U_B \ge 0 ,
\]
which forces $U_B(t_1^-) = 0$ and $U_A(t_1^-) = (k-1)L$. Since $L > 0$, this
gives $N_A \ge 1$ on $I$ --- with $N_A = 0$ we would have $U_A(t_1^-) = 0$ ---
so $U_A$ is strictly decreasing on $I$ and
$U_A(t_1^-) < \sup_I U_A \le (k-1)L$, a contradiction. Exchanging $A$ and $B$
gives the other direction.
\end{proof}

\subsection{The cascade}
\label{app:cascade}

Both tightness results begin from the same construction, which drives the
workload gap of Lemma~\ref{lem:gap} to its ceiling geometrically.

\begin{lemma}
\label{lem:cascade}
Fix $k \ge 2$ and $m \ge 1$ and set $L = k^m$. There is an input and a
work-conserving policy $P$ such that, at time $T_m = mL$, first-come first-served
has completed every job and idles all $k$ servers, while $P$ holds exactly $k-1$
jobs in service, each with remaining service
\[
  f \;=\; \frac{\rho_m}{k-1} \;=\; L\Bigl(1 - \bigl(\tfrac{k-1}{k}\bigr)^{m}\Bigr) ,
  \qquad
  \rho_m \;=\; (k-1)L\Bigl(1 - \bigl(\tfrac{k-1}{k}\bigr)^{m}\Bigr) ,
\]
so that $U_P(T_m) - U_{\mathrm{FCFS}}(T_m) = \rho_m \to (k-1)L$ as
$m \to \infty$. Every $\rho_j$ and every intermediate time is an integer.
\end{lemma}

\begin{proof}
The input consists of $m$ identical rounds; round $j$ occupies
$[T_j, T_{j+1}) = [jL, (j+1)L)$. At time $T_j$ inject, in rank order, $k-1$
\emph{long} jobs of service $L$ and then $L$ \emph{short} jobs of service $1$, so
that the long jobs carry the lower ranks. Let $P$ be the static-priority policy
that dispatches a waiting short job whenever one exists and a long job
otherwise; it is work-conserving.

We show by induction that at $T_j$ first-come first-served has completed
everything and idles, while $P$ has exactly $k-1$ jobs in service with total
remaining $\rho_j$, where $\rho_0 = 0$.

\emph{First-come first-served in round $j$.} All $k$ servers are free at $T_j$,
so it dispatches the $k-1$ long jobs and one short job. Its remaining server
takes the $L$ short jobs back to back and finishes at $T_j + L$; the long jobs
also finish at $T_j + L$. So all $k$ servers are busy on $(T_j, T_{j+1})$ and
the system is empty and idle again at $T_{j+1}$.

\emph{$P$ in round $j$.} At $T_j$, $P$ has one free server, which takes a short
job. Its $k-1$ carried-over jobs all complete at $T_j + \rho_j/(k-1)$ and their
servers take short jobs as well, and $P$ never touches a long job while a short
one waits. So all $k$ servers are busy from $T_j$ until the short jobs run out,
at $T_j + \tau_j$ with $k\tau_j = \rho_j + L$. ($P$ is indeed busy throughout:
the carried-over jobs end at $\rho_j/(k-1) \le \tau_j$, which is
$k\rho_j \le (k-1)(\rho_j + L)$, that is $\rho_j \le (k-1)L$, true by
induction.) At $T_j + \tau_j$ the policy $P$ holds only the $k-1$ long jobs of
the round; it dispatches them and leaves one server idle.

\emph{The gain.} On $(T_j + \tau_j, T_{j+1}]$, first-come first-served has $k$
busy servers and $P$ has $k-1$, so the gap grows at rate $1$ over a window of
length $L - \tau_j$. At $T_{j+1}$ the reference system is empty, and each of
$P$'s $k-1$ long jobs has run for $L - \tau_j$ and has $\tau_j$ left. The
hypothesis is reproduced with $\rho_{j+1} = (k-1)\tau_j = (k-1)(\rho_j + L)/k$,
whose solution with $\rho_0 = 0$ is
$\rho_j = (k-1)L\,(1 - ((k-1)/k)^j)$. With $L = k^m$ every $\rho_j$ and $\tau_j$
for $j \le m$ is an integer.
\end{proof}

\subsection{Tightness of the net-overtake identity}
\label{app:tight1}

\begin{proof}[Proof of Theorem~\ref{thm:tight1}]
\emph{Upward.} Run the cascade of Lemma~\ref{lem:cascade} for $m$ rounds and, at
$T_m$, inject in rank order $k-1$ jobs of service $L$, then the victim $i$ of
service $L$, then $N > 2kL$ jobs of service $1$. Let $P$ continue as the static
priority ``the $k-1$ new long jobs first, then the short jobs, then $i$''.

At $T_m$ all $k$ servers of the reference policy are free and its waiting set in
rank order is the $k-1$ new long jobs, then $i$, then the short ones. It
dispatches the long jobs and then $i$, all in the phase at $T_m$. Hence
$W_{\mathrm{FCFS}}[i] = 0$; the $k-1$ new long jobs are dispatched strictly
before $i$ in its sequence and are in service at $s^{\mathrm{FCFS}}_i = T_m$
with their full service, so $\rho^{\mathrm{FCFS}}_i = (k-1)L$ and
$R^{\mathrm{FCFS}}_i = (k-1)L$.

Under $P$, one server is free at $T_m$ and takes a new long job; when the $k-1$
carried-over jobs finish, all at $T_m + f$ with $f < L$, their servers take the
remaining new long jobs and then short jobs. The victim waits while any short
job waits. Short jobs are executed at rate at most $k$, so they are not exhausted
before $T_m + N/k$, and $N > 2kL > k(f+L)$ makes that later than $T_m + f + L$,
by which time every new long job has completed. So $i$ is dispatched after all
$k-1$ new long jobs have finished, and at $s^P_i$ the other $k-1$ servers hold
short jobs, giving $0 \le \rho^P_i \le k-1$. Also
$R^P_i = \rho_m + (k-1)L$, since the carried-over jobs and the new long jobs all
have rank below $i$ and the new long jobs are untouched at $a_i$. Therefore
$R^P_i - R^{\mathrm{FCFS}}_i = \rho_m$ and
\begin{align*}
  \Delta_i \;&=\; \bigl(R^P_i - R^{\mathrm{FCFS}}_i\bigr) - \rho^P_i
                 + \rho^{\mathrm{FCFS}}_i
  \;\ge\; \rho_m + (k-1)L - (k-1) \\
  \;&=\; 2(k-1)L - (k-1)L\bigl(\tfrac{k-1}{k}\bigr)^{m} - (k-1) ,
\end{align*}
and both error terms are $o(L)$ as $m \to \infty$ with $L = k^m$.

\emph{Downward.} Run the mirrored cascade: inject the $L$ short jobs first in
rank order and the $k-1$ long jobs after, and let $P$ be ``long jobs before short
ones''. The induction of Lemma~\ref{lem:cascade} runs with the roles of the two
policies exchanged, so at $T_m$ the policy $P$ is empty and idle while the
reference holds $k-1$ jobs with $f$ of service left each. At $T_m$ inject, in
rank order, one \emph{filler} job of service $f$, the victim $i$ of service
$L-1$, and $k-1$ jobs of service $L$; let $P$ be ``the $k-1$ new long jobs
first, then $i$, then the filler''.

At $T_m$ the reference policy has $k-1$ servers busy with carried-over jobs,
each with exactly $f$ left, and one free server, which takes the filler --- the
lowest-ranked waiting job --- which also has service $f$. So all $k$ of its
servers complete simultaneously at $T_m + f$, and it then dispatches $i$, with
nothing else dispatched in that phase before it. Hence
$W_{\mathrm{FCFS}}[i] = f$ and $\rho^{\mathrm{FCFS}}_i = 0$.

Under $P$ the system is empty at $T_m$ with $k$ free servers, and it dispatches
the $k-1$ new long jobs and then $i$, all at $T_m$. So $W_P[i] = 0$,
$\mathrm{In}_i = (k-1)L$, $\rho^P_i = (k-1)L$, and the filler, of rank below
$i$, is dispatched after $i$, so $\mathrm{Out}_i = f$. Therefore
\begin{align*}
  \Delta_i \;&=\; k(0 - f) - \bigl((k-1)L - f\bigr)
           \;=\; -(k-1)f - (k-1)L \\
           \;&=\; -\rho_m - (k-1)L \;\longrightarrow\; -2(k-1)L .
\end{align*}

\emph{Non-attainment.} Lemma~\ref{lem:gap} gives
$|R^P_i - R^{\mathrm{FCFS}}_i| < (k-1)L$ strictly, and step (ii) of
Theorem~\ref{thm:identity} gives $0 \le \rho^P_i, \rho^{\mathrm{FCFS}}_i \le
(k-1)L$; combining them in each direction gives $|\Delta_i| < 2(k-1)L$.
\end{proof}

For the record, the exact values of $\Delta_i / L$ produced by these two
families, against the ceiling $2(k-1)$:

\begin{center}
\begin{tabular}{rrrrrr}
\hline
$k$ & $L$ & $m$ & upward $\Delta_i/L$ & downward $\Delta_i/L$ & $2(k-1)$ \\
\hline
2 & 64 & 6 & $127/64 = 1.984$ & $-127/64 = -1.984$ & 2 \\
3 & 81 & 4 & $290/81 = 3.580$ & $-292/81 = -3.605$ & 4 \\
4 & 64 & 3 & $303/64 = 4.734$ & $-303/64 = -4.734$ & 6 \\
\hline
\end{tabular}
\end{center}

In the executed-work convention of Remark~\ref{rem:directions} the upward
column is unchanged to within $o(L)$, while the downward column becomes
$-63/64$, $-130/81$ and $-111/64$, that is $-(k-1)L$ rather than $-2(k-1)L$.
The same-phase share of $|\Delta_i|$ is \devnum{0.00\%}, \devnum{0.69\%} and
\devnum{0.00\%} upward and \devnum{50.39\%}, \devnum{55.48\%} and
\devnum{63.37\%} downward, and the victim's own excess is
\devnum{$+3.00L$}, \devnum{$+3.20L$}, \devnum{$+3.19L$} upward and
\devnum{$-0.98L$}, \devnum{$-0.80L$}, \devnum{$-0.58L$} downward.

\subsection{Tightness of the guard bound}
\label{app:tight2}

\begin{proof}[Proof of Theorem~\ref{thm:tight2}(i)]
At $t = 0$ let there arrive, in rank order, $k$ jobs of service $L$, the victim
$i$, and at least $B + k$ jobs of service $1$; let the base policy be the static
priority ``the $k$ long jobs, then the short ones, then $i$'', and take the
constant budget $B$. First-come first-served dispatches the $k$ long jobs in the
phase at $t = 0$, filling every server, so $W_{\mathrm{FCFS}}[i] = L$.

If $B = 0$ then $E(t)$ contains every waiting job at every dispatch epoch and the
wrapper is first-come first-served, so the excess is $0 = \lceil B/k \rceil$. If
$B > 0$ then at $t = 0$ no job has completed, $E(0)$ is empty, and the base fills
the servers with the $k$ long jobs. They complete at $t = L$, and
$\mathrm{over}[i](L) = 0$ still, every completed job having rank below $i$. From
$t = L$ on, all $k$ servers are free at each integer instant; while
$\mathrm{over}[i] < B$ the set $E$ is empty --- the waiting jobs of rank above
$i$ have no completed job of still higher rank ahead of them --- so the base
dispatches $k$ short jobs, which complete one time unit later. Hence
$\mathrm{over}[i](L + r) = rk$, the first instant at which
$\mathrm{over}[i] \ge B$ is $L + \lceil B/k \rceil$, and there $i$ is the
minimum-rank member of $E$ and is dispatched. So
$\mathrm{excess}_P[i] = \lceil B/k \rceil$.
\end{proof}

\begin{proof}[Proof of Theorem~\ref{thm:tight2}(ii)]
The upper bound $c(k) < 3k-2$ is Theorem~\ref{thm:guard}: for a constant budget,
$k\,\mathrm{excess}_P[i] \le \mathrm{In}_i - \mathrm{Out}_i + 2(k-1)L
 \le \mathrm{In}_i + 2(k-1)L < B + kL + 2(k-1)L = B + (3k-2)L$, the last step
being the strict inequality \eqref{eq:guardin}. When $i$ has no overtaker the
excess is at most $(2-2/k)L$ and the inequality is strict as well.

For the lower bound, fix $k \ge 2$, $m \ge 1$, $L = k^m$, and let $f$ and
$\rho_m$ be as in Lemma~\ref{lem:cascade}. Run the cascade for $m$ rounds. At
$T_m$ inject, all arriving at $T_m$ and in rank order,
\begin{itemize}
  \item $k-1$ \emph{new long} jobs of service $L$, of rank below the victim;
  \item the victim $i$, of service $L$;
  \item one \emph{sync} job of service $f$;
  \item one \emph{pre} job of service $L$;
  \item $k$ \emph{final} jobs of service $L$;
\end{itemize}
set the constant budget to $B = f + L + 1$, and let the base policy be the static
priority ``short jobs of the cascade, then sync, then the new long jobs and pre,
then the finals, and the victim last''.

\emph{The reference run.} At $T_m$ it is empty with all $k$ servers free; it
dispatches the $k-1$ new long jobs and then $i$ in the phase at $T_m$, so
$W_{\mathrm{FCFS}}[i] = 0$.

\emph{The guard does not fire during the cascade.} Inside round $j$ the only
waiting jobs with completed higher-ranked work are that round's long jobs, whose
higher-ranked jobs are the $L$ short jobs of the round and everything in later
rounds, which have not arrived. Over $[T_j, T_j + \tau_j]$ the $k$ servers
execute $k\tau_j = \rho_j + L$ units of work, of which $\rho_j$ belongs to the
carried-over jobs, so at most $L$ units of short-job work can have completed.
Hence $\mathrm{over}[q](t) \le L < B$ for every waiting $q$ at every dispatch
epoch of the cascade, $E(t)$ is empty throughout, and the wrapped run reproduces
Lemma~\ref{lem:cascade}.

\emph{The finale.} At $T_m$ one server is free and $\mathrm{over}[q](T_m) = 0$
for every waiting $q$, so the base chooses and takes the sync job. The $k-1$
carried-over jobs have exactly $f$ left and the sync job has service $f$, so all
$k$ servers free simultaneously at $T_m + f$. There $\mathrm{over}[i] = f$,
because only the sync job has rank above $i$ among the completed jobs, and
$f < B$; no other waiting job has a completed higher-ranked job either, so $E$ is
empty and the base fills all $k$ servers with the $k-1$ new long jobs and pre.
These all have service $L$ and complete simultaneously at $T_m + f + L$, where
\[
  \mathrm{over}[i] \;=\; f + L \;=\; B - 1 \;<\; B ,
\]
since only sync and pre have rank above $i$ among the completed jobs. So $E$ is
still empty, the base fills all $k$ servers with the $k$ final jobs, and they
complete at $T_m + f + 2L$, at which instant
$\mathrm{over}[i] = B - 1 + kL \ge B$. The guard fires, $i$ is the minimum-rank
member of $E$, and it is dispatched with every server free, so $\rho^P_i = 0$.

\emph{The value.} $W_P[i] = f + 2L$, $\mathrm{In}_i = B - 1 + kL$,
$\mathrm{Out}_i = 0$, and
\[
  \frac{k\,\mathrm{excess}_P[i] - B}{L}
  \;=\; \frac{k(f + 2L) - (f + L + 1)}{L}
  \;=\; \frac{(k-1)f + (2k-1)L - 1}{L} .
\]
As $m \to \infty$ with $L = k^m$ we have $f/L \to 1$ and $1/L \to 0$, so the
ratio tends to $(k-1) + (2k-1) = 3k-2$.

The same instance also has $\Delta_i = \rho_m + (k-1)L \to 2(k-1)L$, so the two
sources of slack in Theorem~\ref{thm:guard} are saturated together.
\end{proof}

The values produced, against the ceiling $3k-2$:

\begin{center}
\begin{tabular}{rrrrrrr}
\hline
$k$ & $L$ & $m$ & jobs & $B$ & excess & $(k\,\mathrm{excess}-B)/L$ \\
\hline
2 & 64 & 6 & 396 & 128 & 191 & $127/32 = 3.969$ \\
2 & 128 & 7 & 909 & 256 & 383 & $255/64 = 3.984$ \\
3 & 81 & 4 & 340 & 147 & 227 & $178/27 = 6.593$ \\
3 & 243 & 5 & 1233 & 455 & 697 & $1636/243 = 6.733$ \\
4 & 64 & 3 & 211 & 102 & 165 & $279/32 = 8.719$ \\
4 & 256 & 4 & 1046 & 432 & 687 & $579/64 = 9.047$ \\
5 & 125 & 3 & 399 & 187 & 311 & $1368/125 = 10.944$ \\
6 & 216 & 3 & 677 & 308 & 523 & $1415/108 = 13.102$ \\
\hline
\end{tabular}
\end{center}

\subsection{Consistency, and why the converse fails}
\label{app:consistency}

\begin{proof}[Proof of Proposition~\ref{prop:consistency}]
Induct on dispatch epochs: up to the first epoch at which the two runs could
differ their completion histories agree, hence so does $\mathrm{over}[\cdot]$,
so $E(t)$ is empty in the wrapped run and it copies $A$. The first consequence
uses $\mathrm{over}[q](t) \le \mathrm{In}^A_q$, a completed overtaker of $q$
having been dispatched while $q$ waited; the second uses
Corollary~\ref{cor:necessary}.
\end{proof}

The converse fails, so this is not an equivalence. Take $k = 1$, two jobs of
service $1$ at $t = 0$, the base policy ``higher rank first'', and $B = 1$. Both
runs dispatch in the order $(1,0)$, yet at the second dispatch the waiting job
has $\mathrm{over} = 1 \ge B$, so the guard fires and happens to choose what the
base would have chosen. Nor is there multiplicative consistency: with $k = 1$,
one job of service $L$ followed by $m$ jobs of service $1$ all at $t = 0$, the
wrapper with $B = 0$ is first-come first-served while the base ordering by
service time is not, and the ratio of mean waits is $(2L + m - 1)/(m+1)$,
unbounded in $L$.

\subsection{Approaching the single-server bound}
\label{app:floor}

\begin{proof}[Proof of the second claim of Proposition~\ref{prop:floor}]
Fix $L$, a budget $B$, and a granularity $g$ dividing $B$ with $0 < g \le L$.
Take $k = 1$ and, at $t = 0$, the victim $i$ of rank $0$, then $B/g - 1$ jobs of
service $g$, then one job of service $L$; let the base policy dispatch the jobs
of service $g$ first, then the one of service $L$, and never the victim.
First-come first-served starts $i$ at once, so $W_{\mathrm{FCFS}}[i] = 0$.

Under the wrapper, after $r$ of the jobs of service $g$ have completed,
$\mathrm{over}[i] = rg$, which stays below $B$ for $r \le B/g - 1$. At the last
such instant $\mathrm{over}[i] = B - g < B$, the guard does not fire, and the
base dispatches the job of service $L$. When it completes,
$\mathrm{over}[i] = B - g + L \ge B$ and the guard releases $i$. Hence
$\mathrm{excess}_P[i] = B - g + L$, so $\mathrm{excess}_P[i] - B = L - g$, which
approaches $L$ as $g \to 0$ and never reaches it. For example
$(L, B, g) = (100, 40, 1)$ gives $\mathrm{excess}_P[i] - B = 99$, and
$(64, 32, 1)$ gives $63$.
\end{proof}

\subsection{The price of the guarantee}
\label{app:price}

\begin{proof}[Proof of Theorem~\ref{thm:price}]
Let $S(Q)$ be the set of inverted pairs of a policy $Q$, that is the pairs
$(u,v)$ with $u < v$ in rank and $v \prec u$. Counting the same pairs in two ways,
\[
  \sum_i \mathrm{In}_i \;=\; \sum_{(u,v) \in S(Q)} C_v ,
  \qquad
  \sum_i \mathrm{Out}_i \;=\; \sum_{(u,v) \in S(Q)} C_u ,
\]
exactly, for every $k$. Write $\Delta(Q) = \sum_i (W_{\mathrm{FCFS}}[i] -
W_Q[i])$ for the total saving against the reference and $E = 2(1 - 1/k)L(m+n)$.
Summing Theorem~\ref{thm:identity} over the $m+n$ jobs gives
\begin{equation}
  \Bigl| \Delta(Q) - \tfrac{1}{k} \sum_{(u,v) \in S(Q)} (C_u - C_v) \Bigr|
  \;\le\; E
  \qquad \text{for every } Q .
  \label{eq:sumbound}
\end{equation}

\emph{Numerator.} A pair contributes positively only when $u$ is a long job and
$v$ a short one, contributing $L - s$. For a fixed long job $u$, the number of
short jobs dispatched before it is at most $\mathrm{In}_u / s$, and
Corollary~\ref{cor:necessary} gives $\mathrm{In}_u \le kG + (3k-2)L$. Summing
over the $m$ long jobs and applying \eqref{eq:sumbound},
\[
  \Delta(P) \;\le\; \tfrac{1}{k}\, m\,(L-s)\,\frac{kG + (3k-2)L}{s} \;+\; E .
\]

\emph{Denominator.} Sorting by service time inverts every one of the $mn$
long--short pairs, so its pair sum is exactly $mn(L-s)$, and
\eqref{eq:sumbound} applied to it gives
$\Delta(\mathrm{SJF}) \ge \tfrac{1}{k} m n (L-s) - E$. At $k > 1$ the
denominator carries the same error as the numerator, so treating the reference
saving as exact would understate $\varepsilon$.

\emph{Division.} Both bounds are positive once $mn(L-s)/k > E$, and then
\[
  \frac{\Delta(P)}{\Delta(\mathrm{SJF})}
  \;\le\; \frac{\tfrac{1}{k} m (L-s)\,(kG + (3k-2)L)/s + E}
               {\tfrac{1}{k} m n (L-s) - E}
  \;=\; \frac{kG + (3k-2)L}{n s} \;+\; \varepsilon ,
\]
where $\varepsilon$ collects both corrections. The leading terms are
$\Theta(mn)$ and $E = \Theta(m+n)$, so $\varepsilon = O((m+n)/(mn)) \to 0$.
Dividing numerator and denominator by $m+n$ turns the two total savings into the
two mean savings and leaves the ratio unchanged.
\end{proof}

At $k = 1$ the error term vanishes identically, the sum identity is exact, and
the achievable fraction is $\min(1, G/(ns))$: the only way to beat first-come
first-served in mean wait is to invert pairs in which the later-ranked job is
the shorter one, and the saving is exactly the sum of the service-time
differences over the inversions.
}{}

% ---------------------------------------------------------------------------
% Back matter.  Under the MDPI class each block below is one of the template's
% own macros -- \authorcontributions{}, \funding{}, \institutionalreview{},
% \informedconsent{}, \dataavailability{}, \conflictsofinterest{} -- and the
% headings here are the stand-in for them; main.tex carries the same text in
% those macros.  Anything the authors alone can supply is marked
% [to be completed by the authors] and is not invented.
\section*{Back Matter}
\addcontentsline{toc}{section}{Back Matter}

\paragraph{Author Contributions.}
[to be completed by the authors]. The MDPI template expects the CRediT roles
(conceptualization, methodology, software, validation, formal analysis,
investigation, data curation, writing---original draft preparation,
writing---review and editing, visualization, supervision, project
administration, funding acquisition), each attributed by author initials,
followed by ``All authors have read and agreed to the published version of the
manuscript.''

\paragraph{Funding.}
[to be completed by the authors]. If there is no external funding, the template
asks for ``This research received no external funding''; otherwise the grant
number and the funding body go here.

\paragraph{Institutional Review Board Statement.}
Not applicable.

\paragraph{Informed Consent Statement.}
Not applicable.

\paragraph{Data Availability Statement.}
Five of the six workloads used here are public at the locations cited, and we
redistribute none of them. The CodeBench judging logs are published by the
Universidade Federal do Amazonas \cite{codebench_dataset}; the ACcoding dataset
is published with its data descriptor \cite{accoding2024}; the Open University
Learning Analytics dataset is published with its own \cite{oulad2017}; the Azure
Functions 2021 invocation trace is published in the Azure Public Dataset
repository \cite{azure2021trace}, which asks that \cite{zhang2021harvested} be
cited; and the Intel Netbatch logs are published in the Parallel Workloads
Archive \cite{netbatch2012trace}, which asks that \cite{shai2014netbatch} be
cited and that the logs be acknowledged to their depositors. The sixth, the
continuous-integration trace of the two Firefox hardware worker pools, was
collected by us from Mozilla's public, unauthenticated Taskcluster and
Treeherder APIs \cite{taskcluster_docs,treeherder_docs}; we do not redistribute
it, and the collection scripts that rebuild it from those APIs are part of the
released code. No data were obtained under licence or under an access agreement,
and no personal data were processed beyond the pseudonymous user identifiers the
published judging datasets already carry. The simulator, the predictors, the
guard implementation and the configuration files that reproduce every table and
figure in this paper are to be released with the paper at
[repository URL to be completed by the authors].

\paragraph{Conflicts of Interest.}
[to be completed by the authors]. The MDPI template asks for ``The authors
declare no conflicts of interest'' where that holds, and otherwise for the
competing interest and for a statement on the role of any funder in the design,
collection, analysis or interpretation, in the writing, or in the decision to
publish.
% ---------------------------------------------------------------------------

\bibliography{../docs/related_work/refs}

\end{document}
